\chapter{Entropic Singularities and Topology}
\label{ch:entropic-topology}

The analysis of Chapter~\ref{ch:global} shows that failure of classical
regularity is driven by blow--up of the entropy gradient $\nabla S$.
In this chapter we interpret such singularities as derived geometric
defects in the lamphrodynamic plenum, identify topological invariants
that classify different singularity types, and describe mechanisms by
which derived geometry ``resolves’’ singularities through higher
homotopical structure.
This constitutes the geometric completion of the analytic story.



\section{Derived Singularities of the Plenum Stack}
\label{sec:derived-sing}

Let
\[
\mathrm{Plenum}
= \mathcal{U}
\]
be the derived Artin stack of lamphrodynamic configurations constructed
in Chapter~\ref{ch:plenum-construction}.  By
Theorem~\ref{thm:Main-Plenum}, $\mathrm{Plenum}$ carries a
$2$--shifted symplectic structure whose AKSZ descent produces the
classical BV theory of Chapter~7.

\begin{definition}[Derived Entropic Singularity]
A point $x\in M$ is an \emph{entropic singularity} of a classical
configuration if
\[
|\nabla S(x)|=+\infty.
\]
In the derived picture, the corresponding point in
$\mathrm{Map}(M,\mathrm{Plenum})$ lies outside the classical truncation
but admits a well--defined point in the derived mapping stack.
\end{definition}

\begin{remark}
Derived geometry provides an intrinsic meaning to singular field
configurations: they exist as derived points even when no classical
representative exists.  This is why lamphrodynamics “continues’’
through singularities (cf.~Conjecture~\ref{conj:weak-global}).
\end{remark}



\section{Classification of Singularities}
\label{sec:classification}

Let $U=(\Phi,\vec v,S)$ be classical away from a closed set
$\Sigma_{\mathrm{sing}}\subset M$.
The induced map
\[
U\colon M\setminus\Sigma_{\mathrm{sing}}\longrightarrow \mathrm{Plenum}
\]
extends to a derived morphism
$\widetilde U\colon M\to\mathrm{Plenum}$.
Singularities thus correspond to nontrivial derived fibers at points of
$\Sigma_{\mathrm{sing}}$.

\begin{proposition}[Topological Classification]
\label{prop:classification}
Connected components of $\Sigma_{\mathrm{sing}}$ are classified up to
homotopy by elements of $\pi_{*}(\mathrm{Plenum})$, i.e.\ by homotopy
groups of the lamphrodynamic plenum.
\end{proposition}

\begin{proof}[Proof (Sketch)]
The derived extension $\widetilde U$ replaces the omission of
$\Sigma_{\mathrm{sing}}$ by higher homotopical data in
$\mathrm{Plenum}$, so homotopy classes of defects correspond to
homotopy classes of maps into the target.
\end{proof}

\begin{remark}
This phenomenon parallels topological defects in gauge theory (monopoles,
vortices), but here the relevant target is not a Lie group but the
derived stack $\mathrm{Plenum}$, producing a richer stratification of
defect types.
\end{remark}



\section{Characteristic Classes and Entropy}
\label{sec:char-classes}

The entropy scalar $S$ determines a derived $1$--form $dS$; its
singular behavior influences cohomology classes on $M$.

\begin{definition}[Entropic Characteristic Class]
Define
\[
\mathfrak{c}_{S}:=[dS]\in H^{1}(M\setminus\Sigma_{\mathrm{sing}};\mathbb{R}),
\]
viewed as a de~Rham cohomology class on the nonsingular locus.
\end{definition}

When $\Sigma_{\mathrm{sing}}$ has codimension two or higher, $\mathfrak c_{S}$
extends to a \emph{relative} cohomology class
\[
\mathfrak{c}_{S}\in H^{1}(M,\Sigma_{\mathrm{sing}};\mathbb{R}).
\]

\begin{proposition}[Flux Quantization (conditional)]
\label{prop:flux-quant}
If $\Sigma_{\mathrm{sing}}$ consists of isolated points and
the derived obstruction space has integer lattice structure
(e.g.\ in integral models of entropy),
then the entropic flux through small spheres linking each singular
point is quantized.
\end{proposition}

\begin{proof}[Proof (Sketch)]
Reduction to integral cohomology of the relative pair
$(M,\Sigma_{\mathrm{sing}})$; quantization then follows from integer
obstruction classes.
\end{proof}

\begin{remark}
This resembles magnetic charge quantization but arises purely from
entropy geometry.
\end{remark}



\section{Derived Resolution of Entropic Defects}
\label{sec:resolution}

Given a derived stack $\mathcal X$ and a morphism
$f\colon M\to\mathcal X$ with classical singularities, derived geometry
often admits a resolution---a lift to a homotopy equivalent object
without classical defects.

\begin{theorem}[Derived Resolution]
\label{thm:resolution}
Let $\Sigma_{\mathrm{sing}}$ be a closed set of entropic
singularities of $U$.
Then there exists a derived lift
\[
\widetilde U\colon M\longrightarrow\mathrm{Plenum}
\]
which is a derived resolution of $U$ in the sense that
$\widetilde U$ is a homotopy equivalence on the nonsingular locus and
extends uniquely across $\Sigma_{\mathrm{sing}}$ in the derived
category.
\end{theorem}

\begin{proof}[Proof (Sketch)]
This is a direct consequence of how derived mapping stacks encode
obstruction theory: the obstruction to extending $U$ lives in the
cotangent complex of $\mathrm{Plenum}$, which, by Theorem~\ref{thm:Main-Plenum},
has amplitude $[-1,1]$, forcing obstructions to vanish in strictly
positive degrees.
\end{proof}



\section{Physical Interpretation}
\label{sec:phys-top}

In the analytic picture, entropic singularities are catastrophes of the
entropy gradient. In the derived picture they are homotopy--theoretic
defects classified by topological invariants.  Derived resolution means
that evolution of lamphrodynamic fields remains well--defined even when
classical smoothness fails.  This is consonant with the BV formalism
(Ch.~7): the ghost/antifield sector automatically “fills in’’ missing
information at singular loci.



\section{Summary and Forward Link}
\label{sec:summary12}

We have shown that entropy singularities correspond to derived defects
in the lamphrodynamic plenum, are classified by homotopy invariants,
admit characteristic classes, and possess derived resolutions.
These ideas form the foundation for the geometric reductions of
Chapter~\ref{ch:symmetries-reduction}.

